\addcontentsline{toc}{chapter}{СПИСОК ИСПОЛЬЗОВАННЫХ ИСТОЧНИКОВ}
% \bibliographystyle{utf8gost705u}  % стилевой файл для оформления по ГОСТу
% \bibliography{biblio}          % имя библиографической базы (bib-файла)
\renewcommand{\bibname}{СПИСОК ИСПОЛЬЗОВАННЫХ ИСТОЧНИКОВ}
\begin{thebibliography}{99}
	\bibitem{kurouz_book}
	Джеймс Ф. Куроуз, Кит В. Росс
	\textit{Компьютерные сети}: 2-е издание. ---
	М.: ЗАО Издательский дом «Питер», 2004 --- 432 с. --- (дата обращения: 5.02.2026).
	
	\bibitem{mike_book}
	Майкл Керриски
	\textit{Linux API}: ---
	М.: Изд-во Питер, 2018. --- 400 с. --- (дата обращения: 5.02.2026).
	
	\bibitem{tanen_book}
	Эндрю Таненбаум
	\textit{Компьютерные сети}: 5-е издание. ---
	СПб.: Изд-во Питер, 2018. --- 592 с. --- (дата обращения: 5.02.2026).
	
	\bibitem{man}
	\textit{Мануал Linux}: 
    URL: \{https://man7.org/linux/man-pages/man7/pthreads.7.html\} (дата обращения: 5.02.2026).
	
\end{thebibliography}
%\nocite{*}
